% ============================================================================
% Section 3: SPI Waveform (Task 3)
% Handout requires: calculation of B, SCK/MOSI capture, measured SCK
% frequency, percentage error, manual decoding of the transmitted byte.
% The waveform is the evidence - decode it from the capture, not the code.
%
% PRE-FILLED before measuring: B, the predictions (Expected columns) and the
% explanations. Still to do after the scope capture (all red):
%   - figures/task3_sck_mosi.png
%   - \TSCKMeas, \fSCKMeas, \fSCKErr in values.tex
%   - the "Measured" column of Table (task3-meas) and the "Measured" row of
%     Table (decode)
%   - the two short sentences marked \fillin below
% ============================================================================
\section{SPI Waveform}
\label{sec:spi-waveform}

\paragraph{Group test byte.}
With $n_1 = \nOne$ (\StudentOneNumber) and $n_2 = \nTwo$ (\StudentTwoNumber),
the last three decimal digits of each student number:
\begin{align}
  n_1 \oplus n_2 &= \text{\nOneHex} \oplus \text{\nTwoHex} = \text{\XorHex}
    \quad \text{(\cref{tab:xor})} \\
  B &= \left[(n_1 \oplus n_2) + \text{0x3D}\right] \bmod 256
     = (\XorDec + 61) \bmod 256 = \SumDec = \text{\BHex}.
\end{align}
The result is neither 0x00 nor 0xFF, so the $\oplus\,$0x5A rule does not apply.
The group test byte is therefore $B = \text{\BHex} = \texttt{\BBin}_2$. The build
confirms this hand calculation with a \reg{\_Static\_assert} in \reg{main.c}.

\begin{table}[htbp]
  \centering
  \caption{Bitwise XOR of $n_1$ and $n_2$.}
  \label{tab:xor}
  \begin{tabular}{@{}llll@{}}
    \toprule
     & \textbf{Decimal} & \textbf{Hex} & \textbf{Binary} \\
    \midrule
    $n_1$            & \nOne   & \nOneHex & \texttt{\nOneBin} \\
    $n_2$            & \nTwo   & \nTwoHex & \texttt{\nTwoBin} \\
    $n_1 \oplus n_2$ & \XorDec & \XorHex  & \texttt{\XorBin} \\
    \bottomrule
  \end{tabular}
\end{table}

\paragraph{Capture.}
With \reg{RUN\_TASK}~$=3$, $B$ is written to \reg{SPI2\_DR} every 100~ms while CS
(PB12) is held high, so the EEPROM ignores the transfer. The data register is
written with an 8-bit access; a 16-bit write would pack two frames and produce
16 clocks \cite[\S28.5.9]{rm0091}. CH1 was connected to SCK (PB13) and CH2 to
MOSI (PB15), with both probes and channels set to $\times$10 and the scope
triggered on the first SCK edge.

\scopefig{task3_sck_mosi.png}{One byte ($B = \text{\BHex}$) transmitted by the
  hardware SPI peripheral: CH1 = SCK (PB13), CH2 = MOSI (PB15). Cursors mark one
  SCK period.}{fig:task3}

\begin{table}[htbp]
  \centering
  \caption{Predicted and measured waveform properties (\cref{fig:task3}).}
  \label{tab:task3-meas}
  \begin{tabular}{@{}lll@{}}
    \toprule
    \textbf{Quantity} & \textbf{Expected} & \textbf{Measured} \\
    \midrule
    SCK period $T_\mathrm{SCK}$       & \TSCKPred~\si{\micro\second} & \TSCKMeas~\si{\micro\second} \\
    SCK frequency $f_\mathrm{SCK} = 1/T_\mathrm{SCK}$ & \fSCKPred~kHz & \fSCKMeas~kHz \\
    Clock idle state                  & low (CPOL = 0)        & \fillin{low / high} \\
    Number of clock cycles            & 8 (8-bit frame)       & \fillin{8} \\
    Sampling edge                     & rising (CPHA = 0)     & \fillin{rising / falling} \\
    Bit order                         & MSB first             & \fillin{MSB / LSB first} \\
    Transmitted value                 & \BHex                 & \fillin{0x..} \\
    \bottomrule
  \end{tabular}
\end{table}

\paragraph{Frequency error.}
\begin{equation}
  \%\,\mathrm{error}
  = \frac{\lvert f_\mathrm{measured} - f_\mathrm{predicted}\rvert}{f_\mathrm{predicted}} \times 100
  = \frac{\lvert \fSCKMeas - \fSCKPred \rvert}{\fSCKPred} \times 100
  = \fSCKErr\,\%.
\end{equation}
SCK is derived directly from the HSI RC oscillator, which is factory-calibrated to
$\pm1\,\%$ at 25~\si{\celsius} and specified between $-1\,\%$ and $+2\,\%$ over
0--55~\si{\celsius} \cite[Tab.~37]{stm32f051ds}. The measured error of
\fSCKErr\,\% is \fillin{within / outside} this tolerance\fillin{, with the remainder
attributable to cursor placement on the scope}.

\paragraph{Manual decoding.}
With CPOL = 0 and CPHA = 0, the first SCK edge is rising and is the first data
capture edge; data is latched on every rising edge \cite[\S28.5.6]{rm0091}. The
master therefore places each bit on MOSI half a clock period before the rising
edge that samples it, changing MOSI on the falling edges; the first bit (b7) is
already on MOSI before the first rising edge \cite[Fig.~277, p.~764]{rm0091}.
This matches the EEPROM, which latches SI on the rising edge of SCK
\cite[p.~5]{eeprom}. MOSI was read at each rising edge in \cref{fig:task3}
(\cref{tab:decode}).

\begin{table}[htbp]
  \centering
  \caption{Bit-by-bit decoding of the MOSI trace, MSB first.}
  \label{tab:decode}
  \begin{tabular}{@{}lcccccccc@{}}
    \toprule
    \textbf{Rising edge \#} & 1 & 2 & 3 & 4 & 5 & 6 & 7 & 8 \\
    \textbf{Bit}            & b7 & b6 & b5 & b4 & b3 & b2 & b1 & b0 \\
    \midrule
    \textbf{Expected ($B$)} & 0 & 1 & 0 & 0 & 0 & 1 & 0 & 0 \\
    \textbf{Measured MOSI}  & \fillin{} & \fillin{} & \fillin{} & \fillin{} & \fillin{} & \fillin{} & \fillin{} & \fillin{} \\
    \bottomrule
  \end{tabular}
\end{table}

The decoded byte is \fillin{0100\,0100$_2$ = 0x44}, which \fillin{equals} the group
test byte $B = \text{\BHex}$. Only bits b6 and b2 are high, so on the scope MOSI
should show exactly two high pulses, in the 2nd and 6th bit periods.
